%% Generated by Sphinx.
\def\sphinxdocclass{report}
\documentclass[letterpaper,10pt,english]{sphinxmanual}
\ifdefined\pdfpxdimen
   \let\sphinxpxdimen\pdfpxdimen\else\newdimen\sphinxpxdimen
\fi \sphinxpxdimen=.75bp\relax

\PassOptionsToPackage{warn}{textcomp}
\usepackage[utf8]{inputenc}
\ifdefined\DeclareUnicodeCharacter
% support both utf8 and utf8x syntaxes
  \ifdefined\DeclareUnicodeCharacterAsOptional
    \def\sphinxDUC#1{\DeclareUnicodeCharacter{"#1}}
  \else
    \let\sphinxDUC\DeclareUnicodeCharacter
  \fi
  \sphinxDUC{00A0}{\nobreakspace}
  \sphinxDUC{2500}{\sphinxunichar{2500}}
  \sphinxDUC{2502}{\sphinxunichar{2502}}
  \sphinxDUC{2514}{\sphinxunichar{2514}}
  \sphinxDUC{251C}{\sphinxunichar{251C}}
  \sphinxDUC{2572}{\textbackslash}
\fi
\usepackage{cmap}
\usepackage[T1]{fontenc}
\usepackage{amsmath,amssymb,amstext}
\usepackage{babel}



\usepackage{times}
\expandafter\ifx\csname T@LGR\endcsname\relax
\else
% LGR was declared as font encoding
  \substitutefont{LGR}{\rmdefault}{cmr}
  \substitutefont{LGR}{\sfdefault}{cmss}
  \substitutefont{LGR}{\ttdefault}{cmtt}
\fi
\expandafter\ifx\csname T@X2\endcsname\relax
  \expandafter\ifx\csname T@T2A\endcsname\relax
  \else
  % T2A was declared as font encoding
    \substitutefont{T2A}{\rmdefault}{cmr}
    \substitutefont{T2A}{\sfdefault}{cmss}
    \substitutefont{T2A}{\ttdefault}{cmtt}
  \fi
\else
% X2 was declared as font encoding
  \substitutefont{X2}{\rmdefault}{cmr}
  \substitutefont{X2}{\sfdefault}{cmss}
  \substitutefont{X2}{\ttdefault}{cmtt}
\fi


\usepackage[Bjarne]{fncychap}
\usepackage{sphinx}

\fvset{fontsize=\small}
\usepackage{geometry}


% Include hyperref last.
\usepackage{hyperref}
% Fix anchor placement for figures with captions.
\usepackage{hypcap}% it must be loaded after hyperref.
% Set up styles of URL: it should be placed after hyperref.
\urlstyle{same}

\addto\captionsenglish{\renewcommand{\contentsname}{Contents:}}

\usepackage{sphinxmessages}
\setcounter{tocdepth}{1}



\title{Gio Vase, Raspberry\sphinxhyphen{}radio}
\date{Jun 29, 2020}
\release{1.8}
\author{Author: Monica Amico}
\newcommand{\sphinxlogo}{\vbox{}}
\renewcommand{\releasename}{Release}
\makeindex
\begin{document}

\pagestyle{empty}
\sphinxmaketitle
\pagestyle{plain}
\sphinxtableofcontents
\pagestyle{normal}
\phantomsection\label{\detokenize{index::doc}}



\chapter{Constant}
\label{\detokenize{index:module-constants}}\label{\detokenize{index:constant}}\index{module@\spxentry{module}!constants@\spxentry{constants}}\index{constants@\spxentry{constants}!module@\spxentry{module}}\index{Operation (class in constants)@\spxentry{Operation}\spxextra{class in constants}}

\begin{fulllineitems}
\phantomsection\label{\detokenize{index:constants.Operation}}\pysiglinewithargsret{\sphinxbfcode{\sphinxupquote{class }}\sphinxcode{\sphinxupquote{constants.}}\sphinxbfcode{\sphinxupquote{Operation}}}{\emph{\DUrole{n}{value}}}{}
Type of Operation Enum
\begin{itemize}
\item {} 
JOINED = 0

\item {} 
REFUSED = 1

\item {} 
PING = 2

\item {} 
HUMIDITY = 3

\item {} 
TEMPERATURE = 4

\item {} 
LIGHT = 5

\item {} 
WATER = 6

\item {} 
SET\_HUMIDITY\_MIN = 7

\item {} 
SET\_HUMIDITY\_MAX = 8

\item {} 
SET\_TEMPERATURE\_MIN = 9

\item {} 
SET\_TEMPERATURE\_MAX = 10

\item {} 
SET\_LIGHT\_MAX = 11

\item {} 
SET\_LIGHT\_MIN = 12

\item {} 
SET\_VASE\_PAUSE\_TIME = 13

\item {} 
SET\_VASE\_SEND\_TIME = 14

\item {} 
SET\_RADIO\_PAUSE\_TIME = 15

\item {} 
SET\_RADIO\_DIED\_PING = 16

\item {} 
CONNECTION = 17

\item {} 
DELETED = 18

\item {} 
SET\_WATERING\_LIGHT = 19

\item {} 
SET\_WATER\_CONTAINER\_SIZE = 20

\item {} 
WATER\_CONTAINER\_STATE = 21

\item {} 
RADIO\_JOIN = 22

\item {} 
RADIO\_TRANSMIT\_POWER = 23

\item {} 
VASE\_TRANSMIT\_POWER = 24

\item {} 
ADD\_EXISTING\_VASE = 26

\end{itemize}

\end{fulllineitems}

\index{WaterContainerState (class in constants)@\spxentry{WaterContainerState}\spxextra{class in constants}}

\begin{fulllineitems}
\phantomsection\label{\detokenize{index:constants.WaterContainerState}}\pysiglinewithargsret{\sphinxbfcode{\sphinxupquote{class }}\sphinxcode{\sphinxupquote{constants.}}\sphinxbfcode{\sphinxupquote{WaterContainerState}}}{\emph{\DUrole{n}{value}}}{}
Water container state enum
\sphinxhyphen{} Empty = 0
\sphinxhyphen{} Full = 1

\end{fulllineitems}



\chapter{Request}
\label{\detokenize{index:module-request}}\label{\detokenize{index:request}}\index{module@\spxentry{module}!request@\spxentry{request}}\index{request@\spxentry{request}!module@\spxentry{module}}\index{request() (in module request)@\spxentry{request()}\spxextra{in module request}}

\begin{fulllineitems}
\phantomsection\label{\detokenize{index:request.request}}\pysiglinewithargsret{\sphinxcode{\sphinxupquote{request.}}\sphinxbfcode{\sphinxupquote{request}}}{}{}
Reads the received request from the dashobard server and adds that to the request queue.

\end{fulllineitems}



\chapter{Threads}
\label{\detokenize{index:module-threads}}\label{\detokenize{index:threads}}\index{module@\spxentry{module}!threads@\spxentry{threads}}\index{threads@\spxentry{threads}!module@\spxentry{module}}\index{Reader (class in threads)@\spxentry{Reader}\spxextra{class in threads}}

\begin{fulllineitems}
\phantomsection\label{\detokenize{index:threads.Reader}}\pysiglinewithargsret{\sphinxbfcode{\sphinxupquote{class }}\sphinxcode{\sphinxupquote{threads.}}\sphinxbfcode{\sphinxupquote{Reader}}}{\emph{\DUrole{n}{nome}}}{}
Thread Reader
\index{run() (threads.Reader method)@\spxentry{run()}\spxextra{threads.Reader method}}

\begin{fulllineitems}
\phantomsection\label{\detokenize{index:threads.Reader.run}}\pysiglinewithargsret{\sphinxbfcode{\sphinxupquote{run}}}{}{}
Method representing the thread’s activity.

It calls the function \sphinxstyleemphasis{reader()}

\end{fulllineitems}


\end{fulllineitems}

\index{Writer (class in threads)@\spxentry{Writer}\spxextra{class in threads}}

\begin{fulllineitems}
\phantomsection\label{\detokenize{index:threads.Writer}}\pysiglinewithargsret{\sphinxbfcode{\sphinxupquote{class }}\sphinxcode{\sphinxupquote{threads.}}\sphinxbfcode{\sphinxupquote{Writer}}}{\emph{\DUrole{n}{nome}}}{}
Thread Writer
\index{run() (threads.Writer method)@\spxentry{run()}\spxextra{threads.Writer method}}

\begin{fulllineitems}
\phantomsection\label{\detokenize{index:threads.Writer.run}}\pysiglinewithargsret{\sphinxbfcode{\sphinxupquote{run}}}{}{}
Method representing the thread’s activity.

It calls the function \sphinxstyleemphasis{writer()}

\end{fulllineitems}


\end{fulllineitems}

\index{read\_serial() (in module threads)@\spxentry{read\_serial()}\spxextra{in module threads}}

\begin{fulllineitems}
\phantomsection\label{\detokenize{index:threads.read_serial}}\pysiglinewithargsret{\sphinxcode{\sphinxupquote{threads.}}\sphinxbfcode{\sphinxupquote{read\_serial}}}{}{}
Read a row from the serial port.
\begin{quote}\begin{description}
\item[{Returns}] \leavevmode
request, serial number, signal, param

\item[{Return type}] \leavevmode
str

\end{description}\end{quote}

\end{fulllineitems}

\index{reader() (in module threads)@\spxentry{reader()}\spxextra{in module threads}}

\begin{fulllineitems}
\phantomsection\label{\detokenize{index:threads.reader}}\pysiglinewithargsret{\sphinxcode{\sphinxupquote{threads.}}\sphinxbfcode{\sphinxupquote{reader}}}{}{}
Called from the Thread Reader, to read request/response from radio\sphinxhyphen{}microbit.
if the request/response received is valid send it to the dashboard,
trow HTTP request.

\end{fulllineitems}

\index{write\_serial() (in module threads)@\spxentry{write\_serial()}\spxextra{in module threads}}

\begin{fulllineitems}
\phantomsection\label{\detokenize{index:threads.write_serial}}\pysiglinewithargsret{\sphinxcode{\sphinxupquote{threads.}}\sphinxbfcode{\sphinxupquote{write\_serial}}}{\emph{\DUrole{n}{data}}}{}
Writes data in the serial port.
\begin{quote}\begin{description}
\item[{Parameters}] \leavevmode
\sphinxstyleliteralstrong{\sphinxupquote{data}} (\sphinxstyleliteralemphasis{\sphinxupquote{str}}) \textendash{} data to be written in the serial port

\item[{Returns}] \leavevmode
len(data)

\item[{Return type}] \leavevmode
int

\end{description}\end{quote}

\end{fulllineitems}

\index{writer() (in module threads)@\spxentry{writer()}\spxextra{in module threads}}

\begin{fulllineitems}
\phantomsection\label{\detokenize{index:threads.writer}}\pysiglinewithargsret{\sphinxcode{\sphinxupquote{threads.}}\sphinxbfcode{\sphinxupquote{writer}}}{}{}
Called from the Thread Writer, to take the requests from the queue (dashboard’s requests)
and to send them to the microbit, writing the data of the requests into the serial port.

\end{fulllineitems}



\renewcommand{\indexname}{Python Module Index}
\begin{sphinxtheindex}
\let\bigletter\sphinxstyleindexlettergroup
\bigletter{c}
\item\relax\sphinxstyleindexentry{constants}\sphinxstyleindexpageref{index:\detokenize{module-constants}}
\indexspace
\bigletter{r}
\item\relax\sphinxstyleindexentry{request}\sphinxstyleindexpageref{index:\detokenize{module-request}}
\indexspace
\bigletter{t}
\item\relax\sphinxstyleindexentry{threads}\sphinxstyleindexpageref{index:\detokenize{module-threads}}
\end{sphinxtheindex}

\renewcommand{\indexname}{Index}
\printindex
\end{document}